\section{Discussion}
\label{sec:discussion}

\subsection{What the experiment establishes}
Across three source regimes, authentic ISV corpus priming is associated with generally positive but strongly configuration-dependent shifts in canonical coverage---with the necessary caveat that HIGH contains negative Gemini mean $\Delta$s and that priming is not universal.
Matched-six mean $\Delta$ is positive on HIGH, LOW, and UNSEEN; UNSEEN shows uniform replicate positivity.
Qualitative audits show that positive $\Delta$ can occur with or without corpus-opening near-copy, depending on regime (\Cref{sec:qual}).
Operationally, execution integrity (same-model pairing; recorded protocol deviations) belongs in the scientific record because one invalid cell and one multi-message constraint materially affect interpretation.

\subsection{What it does not establish}
The study does not establish that priming always improves translation quality, that overlap causes larger effects, that any model is best, or that observed $\Delta$ is causally produced by corpus content alone.
It also does not establish statistical significance.
Readers should resist converting matched-six means into a claim that ``priming improves ISV by about 6--9 points'' as if that were a portable constant.

\subsection{Why cross-regime consistency and heterogeneity matter}
If positive $\Delta$ appeared only on HIGH, a plausible scepticism would be that results merely reward thematic overlap.
LOW and UNSEEN reduce that scepticism for the coverage metric, while still leaving length/genre confounds.
Cross-regime persistence is therefore evidence of broader association within the tested regimes, not proof of a universal intervention.
At the same time, Gemini's HIGH$\rightarrow$LOW reversal shows that regime-averaged slogans fail: the same priming corpus can yield differently signed outcomes by configuration (\Cref{sec:results,sec:cross}).

\subsection{Metric vs.\ quality}
Because coverage is sensitive to inventory hits and orthography, textual pathways that move the metric are only partially aligned with naive notions of ``better translation.''
A Primed output can gain coverage by adopting inventory-supported forms while still containing naming instability, formatting artifacts, or unassessed adequacy errors.
Future work should add human evaluation protocols before claiming quality gains.

\subsection{A restrained conceptual framing}
A useful descriptive framing---not a formal causal law---is that priming appears to be a context-sensitive intervention whose observed effect depends on the interaction between model configuration and source regime.
``Success'' is therefore regime-conditional and configuration-conditional: HIGH-only positives and LOW/UNSEEN positives need not coincide.

\subsection{Relation to in-context learning literature}
Prior LLM-MT work already shows that prompt content and domain matter~\citep{hendy2023good,moslem2023adaptive,zhu2023multilingual}.
Our contribution measures a specific Direct/Primed contrast for Polish-to-ISV under three frozen source regimes with paired repeats.
The ISV coverage metric makes the result auditable inside this repository even though it is not a substitute for human quality evaluation.

\subsection{External validity}
The study uses contemporary vendor interfaces and named configurations recorded at collection time.
External validity to other interfaces, later model snapshots, or open-weights decoding stacks is unknown.
Internal validity for the reported paired contrasts depends on the frozen sources, frozen corpus hash, and documented invalidation of the mismatched Qwen LOW cell.
